\chapter{Селектор центрального симплекса масс дополнительных рёбер}
\label{ch:v8-baryon-c0-extra-edge-mass-central-trace-simplex-selector}

\section{Нормированный симплекс и две конкурирующие конвенции}

Предыдущий гейт локализовал относительную свободу в трёх положительных
центральных весах. После удаления общего масштаба положим
\begin{equation}
 p_Y+p_u+p_d=1,\qquad p_Y,p_u,p_d>0,
 \label{eq:v8-extra-edge-simplex-normalization}
\end{equation}
так что массовый вектор равен
\begin{equation}
 \boldsymbol\mu(p)=(2p_Y,3p_u,3p_d).
 \label{eq:v8-extra-edge-simplex-mass-map}
\end{equation}
Равная микроскопическая плотность следа на трёх блоках даёт
\begin{equation}
 p^{\rm count}=\left(\frac13,\frac13,\frac13\right),\qquad
 \boldsymbol\mu(p^{\rm count})=\left(\frac23,1,1\right),
 \label{eq:v8-extra-edge-simplex-counting-candidate}
\end{equation}
то есть отношение $2:3:3$. Но столь же положительная конвенция равной массы
на ребро требует
\begin{equation}
 p^{\rm edge}=\left(\frac37,\frac27,\frac27\right),\qquad
 \boldsymbol\mu(p^{\rm edge})=\left(\frac67,\frac67,\frac67\right).
 \label{eq:v8-extra-edge-simplex-equal-edge-candidate}
\end{equation}
Эти два вектора не пропорциональны. Первая конвенция считает одинаково
каждую внутреннюю компоненту, вторая --- каждый edge-канал; текущая
архитектура не содержит правила выбора между ними.

\section{Максимум энтропии как условный near miss}

Real-завершённые размеры трёх блоков равны $(4,6,6)$. Центральная матрица
плотности с собственными значениями $(r_Y,r_u,r_d)$ удовлетворяет
\begin{equation}
 4r_Y+6r_u+6r_d=1,
 \qquad
 S(r)=-4r_Y\log r_Y-6r_u\log r_u-6r_d\log r_d.
 \label{eq:v8-extra-edge-simplex-entropy}
\end{equation}
Единственный максимум равен
\begin{equation}
 r_Y=r_u=r_d=\frac1{16},\qquad
 \Hess S=\operatorname{diag}(-64,-96,-96).
 \label{eq:v8-extra-edge-simplex-entropy-maximum}
\end{equation}
Следовательно, максимум энтропии воспроизводит равную микроскопическую
плотность и кандидат $2:3:3$. Но энтропийный функционал не является членом
текущего родительского действия; как и в прежнем singlet--triplet-аудите,
это дополнительный принцип, а не происхождение trace-state.

\section{Симметрия и унимодулярность}

Три представления $(1,2)_{1/2}$, $(3,1)_{5/3}$ и $(3,1)_{2/3}$ попарно
неэквивалентны. Поэтому сохраняющая калибровочные типы группа действует на
центре тривиально:
\begin{equation}
 \dim Z(\mathcal A_z)^G=3.
 \label{eq:v8-extra-edge-simplex-fixed-center}
\end{equation}
Нет допустимой перестановки, которая отождествляла бы все три веса.

Унимодулярность полного блока ограничивает фазовые генераторы условием
\begin{equation}
 4\theta_Y+6\theta_u+6\theta_d=0.
 \label{eq:v8-extra-edge-simplex-unimodularity}
\end{equation}
Эта строка имеет ранг один и оставляет двумерное фазовое ядро. Она не
содержит $p_Y,p_u,p_d$ и потому не уменьшает положительный trace-симплекс.
Такое разделение согласуется с ролью унимодулярности как ограничения
калибровочной группы и аномалий, а не как выбора состояния
\cite{v8-extra-edge-simplex-unimodularity-source}.

\section{Gauge matching и KMS являются калибровками}

Для матрицы индексов $\mathcal I$ предыдущего гейта
\begin{equation}
 \mathcal I p^{\rm count}
 =\left(\frac{61}{18},\frac16,\frac13\right),\qquad
 \mathcal I p^{\rm edge}
 =\left(\frac{125}{42},\frac3{14},\frac27\right).
 \label{eq:v8-extra-edge-simplex-gauge-images}
\end{equation}
Поскольку $\det\mathcal I=-7/4$, заранее заданная тройка gauge-коэффициентов
однозначно восстанавливает $p$. Но без независимой слепой мишени это
обратная калибровка, а не внутреннее предсказание spectral action
\cite{v8-extra-edge-simplex-spectral-action}.

Тепловой класс также реализует весь симплекс. Для центрального
гамильтониана
\begin{equation}
 h=\varepsilon_YP_Y+\varepsilon_uP_u+\varepsilon_dP_d
 \label{eq:v8-extra-edge-simplex-central-hamiltonian}
\end{equation}
отношения весов имеют вид
\begin{equation}
 \frac{p_u}{p_Y}=e^{-\beta(\varepsilon_u-\varepsilon_Y)},\qquad
 \frac{p_d}{p_Y}=e^{-\beta(\varepsilon_d-\varepsilon_Y)}.
 \label{eq:v8-extra-edge-simplex-gibbs-parameterization}
\end{equation}
Например, равномассовая точка требует двух равных разрывов
\begin{equation}
 \varepsilon_u-\varepsilon_Y
 =\varepsilon_d-\varepsilon_Y
 =\frac1\beta\log\frac32.
 \label{eq:v8-extra-edge-simplex-equal-mass-gaps}
\end{equation}
Подробный баланс связывает скорости после задания этих разрывов, но не
порождает их \cite{v8-extra-edge-simplex-detailed-balance}.

\section{Стационарность trace-весов}

Если формально варьировать веса в квадратичном действии, получается
\begin{equation}
 S_z(p,z)=2p_Y|z_Y|^2+3p_u|z_u|^2+3p_d|z_d|^2.
 \label{eq:v8-extra-edge-simplex-weight-action}
\end{equation}
На физическом вакууме $z=0$ выполнено
\begin{equation}
 \nabla_pS_z\big|_{z=0}=0,\qquad
 \Hess_pS_z\big|_{z=0}=0_{3\times3}.
 \label{eq:v8-extra-edge-simplex-vacuum-flatness}
\end{equation}
Поэтому вакуумная стационарность оставляет весь симплекс плоским. Вне
вакуума внутренняя стационарная точка возможна лишь при равенстве трёх
коэффициентов $2|z_Y|^2=3|z_u|^2=3|z_d|^2$; иначе линейный функционал
стремится к границе и разрушает верность состояния.

Итоговый реестр равен
\begin{equation}
 N_{\rm conditional\ conventions}=2/2,\qquad
 N_{\rm intrinsic\ selectors}=0/8,
 \label{eq:v8-extra-edge-simplex-ledgers}
\end{equation}
где проверены счётный след, равная edge-масса, максимум энтропии,
type-симметрия, унимодулярность, gauge matching, KMS/Gibbs и вакуумная
стационарность.

\begin{theorem}[No-go внутреннего селектора trace-симплекса]
Счётный след и максимум энтропии условно выбирают отношение $2:3:3$, тогда
как равная масса на ребро выбирает $1:1:1$. Все три текущие симметрии
сохраняют трёхмерный центр; унимодулярность ограничивает только фазы,
gauge matching требует внешней мишени, KMS параметризует симплекс двумя
разрывами, а вакуумная стационарность полностью плоска. Поэтому внутренний
селектор точки симплекса отсутствует.
\end{theorem}

\begin{nogocriterion}[Запрет энтропийного или gauge-постдикта]
Нельзя считать $2:3:3$ выведенным лишь потому, что оно совпадает с
максимально смешанным состоянием. Нельзя также подставлять наблюдаемые
gauge-коэффициенты в обратимую матрицу индексов и называть полученный
trace-state слепым предсказанием. Оба маршрута требуют нового входа.
\end{nogocriterion}

\begin{protocol}\begin{enumerate}
\item размерность нормированного симплекса: \textbf{$2$};
\item счётный trace-кандидат: \textbf{$2:3:3$};
\item равная edge-масса: \textbf{$1:1:1$};
\item максимум энтропии: \textbf{равная микроскопическая плотность};
\item энтропийный принцип унаследован: \textbf{нет};
\item type-сохраняющий неподвижный центр: \textbf{размерность $3$};
\item унимодулярность ограничивает веса: \textbf{нет};
\item gauge matching без внешней мишени: \textbf{не выбирает};
\item независимых Gibbs-разрывов: \textbf{два};
\item вакуумный гессиан по весам: \textbf{ранг $0$};
\item условных конвенций: \textbf{$2/2$};
\item внутренних селекторов: \textbf{$0/8$};
\item минимальных новых относительных параметров: \textbf{два};
\item trace-симплекс закрыт: \textbf{нет};
\item следующий гейт: \textbf{минимальные данные центрального гамильтониана}.
\end{enumerate}\end{protocol}

Машинный сертификат сохранён в
\path{s2t_v8_baryon_c0_singlet_triplet_central_gap_isotypic_channel_extra_edge_mass_central_trace_simplex_selector_gate_results.json}.

\begin{thebibliography}{9}
\bibitem{v8-extra-edge-simplex-spectral-action}
A.~H.~Chamseddine, A.~Connes,
\emph{The Spectral Action Principle}, arXiv:hep-th/9606001.
\bibitem{v8-extra-edge-simplex-unimodularity-source}
E.~Alvarez, J.~M.~Gracia-Bond\'ia, C.~P.~Mart\'in,
\emph{Anomaly Cancellation and Gauge Group of the Standard Model in NCG},
arXiv:hep-th/9506115.
\bibitem{v8-extra-edge-simplex-detailed-balance}
F.~Fagnola, V.~Umanit\`a,
\emph{Generators of Detailed Balance Quantum Markov Semigroups},
arXiv:0707.2147.
\end{thebibliography}